% Préambule commun aux manuels de Martinet (français et anglais).
% Compilation : pdflatex manuel-fr.tex (deux fois), idem manuel-en.tex, depuis docs/src/.
\documentclass[11pt,a4paper]{article}
\usepackage[utf8]{inputenc}
\usepackage[T1]{fontenc}
\usepackage{lmodern}
\usepackage{microtype}
\usepackage[a4paper,margin=2cm,headheight=24pt]{geometry}
\usepackage{graphicx}
\usepackage[table]{xcolor}
\usepackage[most]{tcolorbox}
\usepackage{enumitem}
\usepackage{booktabs}
\usepackage{array}
\usepackage{tabularx}
\usepackage[nobottomtitles*]{titlesec}
\usepackage{needspace}
\usepackage{fancyhdr}
\usepackage{fontawesome5}
\usepackage{etoolbox}
\DeclareUnicodeCharacter{2192}{\ensuremath{\rightarrow}}
\DeclareUnicodeCharacter{2191}{\ensuremath{\uparrow}}
\DeclareUnicodeCharacter{2193}{\ensuremath{\downarrow}}
\DeclareUnicodeCharacter{25B2}{\ensuremath{\blacktriangle}}
\DeclareUnicodeCharacter{25BC}{\ensuremath{\blacktriangledown}}
\DeclareUnicodeCharacter{21BA}{\ensuremath{\circlearrowleft}}
\DeclareUnicodeCharacter{21BB}{\ensuremath{\circlearrowright}}
\usepackage{amssymb}

% Palette de l'application : bleu ciel, rouge brique pour les alertes
\definecolor{bordeaux}{HTML}{14548C}
\definecolor{bordeauxl}{HTML}{E6F0F8}
\definecolor{alert}{HTML}{B5382B}
\definecolor{alertl}{HTML}{F8E6E3}
\definecolor{or}{HTML}{B8860B}
\definecolor{orl}{HTML}{FBF3DC}
\definecolor{sapin}{HTML}{2C7A55}
\definecolor{sapinl}{HTML}{E2F1E9}
\definecolor{sky}{HTML}{1C6DB4}
\definecolor{skyl}{HTML}{E6F0F8}
\definecolor{summit}{HTML}{6B7A8F}
\definecolor{night}{HTML}{2A2420}
\definecolor{ink}{HTML}{1F2430}
\AtBeginDocument{\color{ink}}

\titleformat{\section}
  {\Large\bfseries\sffamily\color{bordeaux}}
  {\makebox[1.6em][l]{\thesection}}{0.3em}{}
  [{\vspace{1pt}\color{bordeaux!35}\titlerule[1.3pt]}]
\titlespacing*{\section}{0pt}{14pt}{7pt}
\titleformat{\subsection}{\large\bfseries\sffamily\color{night}}{\thesubsection}{0.5em}{}
\pretocmd{\section}{\needspace{10\baselineskip}}{}{}
\pretocmd{\subsection}{\needspace{7\baselineskip}}{}{}
\widowpenalty=10000
\clubpenalty=10000
\raggedbottom
\setlength{\parindent}{0pt}
\setlength{\parskip}{6pt}
\setlist[enumerate]{itemsep=3pt,leftmargin=6mm}
\setlist[itemize]{itemsep=2pt,leftmargin=5mm}

\pagestyle{fancy}
\fancyhf{}
\renewcommand{\headrulewidth}{0.4pt}
\renewcommand{\headrule}{\hbox to\headwidth{\color{bordeaux!40}\leaders\hrule height \headrulewidth\hfill}}
\fancyfoot[C]{\sffamily\small\color{summit}\thepage}
\renewcommand{\sectionmark}[1]{\markboth{#1}{}}

\newcommand{\menu}[1]{\textsf{\textbf{\textcolor{bordeaux}{#1}}}}
\newcommand{\touche}[1]{\tcbox[on line,boxsep=0pt,left=2pt,right=2pt,top=1.5pt,bottom=1.5pt,
  colback=white,colframe=summit,boxrule=0.5pt,arc=1.5pt]{\footnotesize\textsf{#1}}}
\newcommand{\fichier}[1]{\texttt{\small #1}}

\newcommand{\boite}[5]{\newtcolorbox{#1}{enhanced,unbreakable,colback=#2,colframe=#3,
  coltitle=white,fonttitle=\bfseries\sffamily\small,boxrule=0.9pt,arc=2.5pt,
  left=2.5mm,right=2.5mm,top=1.2mm,bottom=1.2mm,title={#4~~#5}}}
\newcommand{\capture}[2][1]{\par\needspace{8\baselineskip}\begin{center}\setlength{\fboxsep}{0pt}\setlength{\fboxrule}{0.6pt}%
  \fcolorbox{summit!50}{white}{\includegraphics[width=\dimexpr#1\linewidth-1.3pt\relax]{#2}}\end{center}}

\usepackage[english]{babel}
\usepackage[colorlinks=true,linkcolor=bordeaux,urlcolor=bordeaux]{hyperref}
\hypersetup{pdftitle={Martinet — user guide},pdfauthor={},pdfcreator={Martinet}}
\boite{astucebox}{orl}{or!85!black}{\faLightbulb}{Good to know}
\boite{securitebox}{sapinl}{sapin}{\faShield*}{What Martinet never does}
\boite{attentionbox}{alertl}{alert}{\faExclamationTriangle}{Caution}
\newcommand{\astuce}[1]{\par\needspace{5\baselineskip}\begin{astucebox}#1\end{astucebox}}
\newcommand{\attention}[1]{\par\needspace{5\baselineskip}\begin{attentionbox}#1\end{attentionbox}}
\fancyhead[L]{\sffamily\small\color{bordeaux}\textbf{MARTINET — USER GUIDE}}
\fancyhead[R]{\sffamily\small\color{summit}\leftmark}
\input{version}

\begin{document}

\begin{titlepage}
\centering
\vspace*{2.2cm}
\includegraphics[width=3.6cm]{../../logo_512.png}\par
\vspace{8mm}
{\sffamily\Huge\bfseries\color{bordeaux} MARTINET\par}
\vspace{3mm}
{\sffamily\LARGE\color{ink} Keep Windows up to date in one go\par}
\vspace{2mm}
{\sffamily\large\color{summit} User guide \textbullet{} version \versionmartinet\par}
\vspace{12mm}
\begin{tcolorbox}[enhanced,width=0.82\linewidth,colback=skyl,colframe=sky,arc=4pt,boxrule=1.2pt,left=5mm,right=5mm,top=3mm,bottom=3mm]
\sffamily\centering
The swift (\textit{martinet} in French) spends its life on the wing and hardly ever lands:\\
Martinet does its rounds on its own.\\[2mm]
\raggedright
\textbf{Check everything} \quad → \quad \textbf{Update and clean up} \quad → \quad \textbf{Schedule}
\end{tcolorbox}
\vfill
{\footnotesize\color{summit}\sffamily Windows 10 and 11\par}
\end{titlepage}

\tableofcontents
\newpage

\section{What Martinet is for}

On Linux, one command (\texttt{apt}, \texttt{yay}…) updates everything. On Windows, each program has
its own updater, drivers have another one, Windows yet another. Martinet brings all of it into one
window:

\begin{center}
\renewcommand{\arraystretch}{1.35}
\begin{tabularx}{\linewidth}{>{\raggedright\arraybackslash}X >{\sffamily\bfseries\color{bordeaux}}l c}
\toprule
\textbf{I want to…} & \textbf{\color{ink}Page} & \textbf{Section} \\
\midrule
see at a glance what is out of date & Dashboard & \ref{sec:accueil} \\
update my programs without opening them one by one & Programs (winget) & \ref{sec:programmes} \\
also track programs winget does not know & Outside winget & \ref{sec:hors} \\
install Windows Update, drivers, the graphics driver & Windows and drivers & \ref{sec:windows} \\
find shortcuts that lead nowhere & Dead shortcuts & \ref{sec:raccourcis} \\
free disk space & Clean-up & \ref{sec:menage} \\
have all of this done without me & Schedule & \ref{sec:planif} \\
\bottomrule
\end{tabularx}
\end{center}

\begin{securitebox}
\begin{itemize}
\item It does not “clean” the registry: no measurable gain, a real risk of breaking a program.
\item It deletes no shortcut: they are \textbf{set aside} and restored in one click.
\item It touches no personal file; browser caches are emptied only if you check them.
\item It only downloads an installer from the vendor's \textbf{official} address, over HTTPS.
\item It never queries a network drive: a sleeping NAS stays asleep.
\end{itemize}
\end{securitebox}

\section{Installing and starting}

Download \fichier{Martinet-Setup.exe} from the project's \textbf{Releases} page and run it. It
installs into your account, without administrator rights. Martinet then updates itself: it offers to
at start-up when a new version is published.

At start-up, Windows asks \textbf{once} for administrator permission: Windows Update, part of the
clean-up and scheduling need it. If you decline, Martinet still works; a yellow banner reminds you
what is missing and offers to restart with the rights.

On the very first start, Martinet offers to turn on the \textbf{scheduled run} every Sunday at noon
(section~\ref{sec:planif}).

\section{The dashboard}\label{sec:accueil}

\capture{../capture_accueil-en.png}

\begin{enumerate}
\item \menu{Check everything} queries winget, vendors, Windows Update and NVIDIA, looks for dead
shortcuts and measures the clean-up. \textbf{Nothing is changed} while checking, which takes one to
three minutes (Windows Update is the slowest).
\item Each card sums up a page; click it to open the page.
\item \menu{Update and clean up} does everything that can be done without you: winget programs,
Windows updates (optional ones included if you chose so), Defender signatures, the checked clean-up,
and, according to the Schedule page, the NVIDIA driver and dead shortcuts set aside. Before starting,
Martinet shows the exact list, with the numbers, and what it will not do.
\end{enumerate}

\astuce{\menu{Show details}, bottom right, unfolds the live log: every line written by winget,
Windows Update or DISM. \menu{Cancel} stops the current operation and the waiting ones.}

\section{Programs (winget)}\label{sec:programmes}

\textbf{winget} is the package manager built into Windows 10 and 11: the equivalent of \texttt{yay}.
It knows thousands of programs and the Microsoft Store apps, and updates them silently.

\capture{../capture_programmes-en.png}

\begin{itemize}
\item Checked rows go with \menu{Update selected}, one by one: one failure does not stop the others,
and the summary says which ones failed.
\item \menu{Ignore / stop ignoring} removes a program from updates, manual and scheduled alike (a
program you deliberately keep at an older version, for instance).
\item \menu{Show ignored} brings them back, greyed out.
\end{itemize}

\attention{If winget is missing (older Windows 10), install “App Installer” from the Microsoft Store.}

\section{Outside winget}\label{sec:hors}

Some programs are not tracked by winget: installed by hand from the vendor's site, or unknown to its
catalogue. Martinet sorts them into two lists.

\capture{../capture_hors-en.png}

\subsection{Programs with a known source}

Martinet reads the latest version where the vendor publishes it:
\begin{itemize}
\item the project's \textbf{GitHub releases} (“dev”, “rc”, “beta” pre-releases are skipped);
\item the vendor's \textbf{download page}, read with a pattern.
\end{itemize}
\menu{Download and install} fetches the official installer when its address is known and runs it:
you follow the installation on screen. Otherwise \menu{Open the page} opens the download page.

The shipped catalogue covers PHD2, ASCOM Platform, N.I.N.A., ASTAP, GraXpert and SharpCap.

\subsection{Programs without a known source}

For each one, three options:
\begin{enumerate}
\item \menu{Hand over to winget…} looks for the program in the winget catalogue; if found, winget
reinstalls it over the existing one (settings are usually kept) and tracks it from then on.
\item \menu{Add a source…} tells Martinet where to find its latest version:
\capture[0.62]{../capture_source-en.png}
For a GitHub project, “owner/repository” is enough (the full address is accepted). For a web page,
the version pattern is a regular expression whose first group in brackets is the version; the highest
one found on the page is kept. \menu{Test} shows what Martinet finds before saving.
\item \menu{Ignore} removes it from the list (drivers, tools without updates…).
\end{enumerate}

\astuce{\menu{Sources…} lists the catalogue and your sources. Editing a catalogue source creates your
own version, which replaces it. \menu{Import…} loads a ready-made source file. Programs that update
themselves (browsers, Microsoft 365, Adobe, NVIDIA App…) are marked “Updates itself”.}

\section{Windows and drivers}\label{sec:windows}

\capture{../capture_windows-en.png}

\subsection{Graphics card}

For an \textbf{NVIDIA} card, Martinet queries NVIDIA's official service (the one the NVIDIA app uses)
with the card's identifier and compares with the installed driver. \menu{Download and install the
driver} fetches the official installer and runs it. For an AMD or Intel card, \menu{Open the page}
leads to the vendor's detection tool.

\subsection{Windows Update}

The list is the one in \textit{Settings > Windows Update}, vendor drivers included (chipset, Wi-Fi,
audio…). \textbf{Optional} updates (optional drivers, Windows previews) are included as long as
\menu{Include optional updates} is checked; this choice also applies to “Update and clean up” and to
scheduled runs. \menu{Update Defender} refreshes the antivirus signatures.

\attention{After installing, Windows often asks for a restart: Martinet offers it (in one minute, time
to save your work).}

\section{Dead shortcuts}\label{sec:raccourcis}

\capture{../capture_raccourcis-en.png}

Martinet reads the Desktop, Start menu and taskbar shortcuts and keeps those whose target has
disappeared. To be safe, it \textbf{does not count}:
\begin{itemize}
\item special shortcuts (Control Panel, Store apps, installer shortcuts);
\item targets on the network or on a missing drive (USB stick, unplugged disk).
\end{itemize}
\menu{Set selected aside} moves them into Martinet's data, with their original location. The lower
list lets you \menu{Restore} them at any time.

\section{Clean-up}\label{sec:menage}

\capture{../capture_menage-en.png}

Only what Windows can recreate: temporary files older than 7 days, error reports, the Windows Update
download cache, installers already used, crash memory dumps, Windows logs older than 7 days, driver setup
leftovers (\texttt{C:\textbackslash NVIDIA}, \texttt{C:\textbackslash AMD}). Never checked by default: the
thumbnail cache, browser caches (emptied only with the browsers closed; no history, cookies or passwords),
shader caches (stutter in games while they rebuild), superseded Windows components (DISM, several
minutes) and the \textbf{recycle bin}, which asks for confirmation. A file in use is left in place.

\section{Schedule: the automatic mode}\label{sec:planif}

\capture{../capture_planif-en.png}

Check \menu{Scheduled run}, choose the day and time, then \menu{Save}. Martinet creates a Windows
scheduled task that:
\begin{itemize}
\item runs with your highest rights, only while your session is open;
\item catches up on a missed run at the next start-up (computer switched off);
\item does what is checked under “What is done without you”: programs, Windows Update and drivers,
the NVIDIA driver (installed without a window), dead shortcuts set aside, clean-up;
\item shows a \textbf{notification} summing up the run at the end; click it to open Martinet.
\end{itemize}

Programs outside winget are only \textbf{reported}: not all of their installers can run without you.
\menu{Run it now} does the same work straight away.

\section{History}

\capture{../capture_historique-en.png}

Every run, manual or scheduled, with line-by-line details. \menu{File > Open the data folder} leads to
the folder where Martinet keeps its settings, sources, history and the shortcuts set aside
(\fichier{\%LOCALAPPDATA\%\textbackslash Martinet}).

\section{Command line}

\begin{center}
\begin{tabular}{>{\ttfamily}l l}
\toprule
run.py --verifier & lists what needs updating, without changing anything \\
run.py --auto & full run, like the scheduled task \\
run.py --lang fr & French interface \\
\bottomrule
\end{tabular}
\end{center}

\section{If something goes wrong}

\begin{itemize}
\item \menu{Help > Report a problem…} sends your description with the Martinet and system versions,
nothing else.
\item A failing step (no network, Windows Update service stopped…) does not stop the others: Martinet
says which one, and why.
\item A program that refuses to update through winget is often open: close it and try again.
\end{itemize}

\end{document}
